\documentclass[11pt]{article}
\usepackage{amsmath, amssymb, amsthm}
\usepackage{graphicx, hyperref, geometry, booktabs}
\geometry{margin=1in}

\newtheorem{definition}{Definition}
\newtheorem{theorem}{Theorem}
\newtheorem{lemma}{Lemma}
\newtheorem{corollary}{Corollary}
\newtheorem{proof}{Proof}

\title{A Formal Mathematical Framework for Bias Interaction Effects in LLM-as-a-Judge}
\author{Student A, Student B}
\date{July 2026}

\begin{document}
\maketitle

\begin{abstract}
We present a formal mathematical framework for understanding bias interaction effects in LLM-as-a-Judge systems. We define bias as a function of prompt perturbations, introduce the Interaction Ratio (IR) as a metric for quantifying non-additive bias effects, and prove several properties about bias interactions under different model architectures. Our framework unifies the 35+ documented bias types under a common mathematical language and identifies the conditions under which biases compound, cancel, or remain additive.
\end{abstract}

\section{Definitions}

\begin{definition}[Score Function]
Let $\mathcal{M}$ be an LLM judge model. Let $\mathcal{P}$ be the space of scoring prompts. The score function $S_\mathcal{M}: \mathcal{P} \rightarrow [1,5]$ maps a prompt $p \in \mathcal{P}$ to an integer score in $\{1,2,3,4,5\}$.
\end{definition}

\begin{definition}[Baseline Prompt]
A baseline prompt $p_0 \in \mathcal{P}$ is a scoring prompt with:
\begin{itemize}
    \item Ascending rubric order (Score 1, Score 2, ..., Score 5)
    \item Arabic numeral score IDs $\{1,2,3,4,5\}$
    \item No reference answer
\end{itemize}
The baseline score is $s_0 = S_\mathcal{M}(p_0)$.
\end{definition}

\begin{definition}[Perturbation Operator]
A perturbation operator $\phi: \mathcal{P} \rightarrow \mathcal{P}$ modifies a specific component of the scoring prompt. We define:
\begin{itemize}
    \item $\phi_{ro}$: Perturb rubric order (ascending $\to$ descending/random)
    \item $\phi_{si}$: Perturb score IDs (Arabic $\to$ letter/Roman)
    \item $\phi_{ra}$: Perturb reference answer (add/change reference score)
    \item $\phi_{pos}$: Perturb response position (first/second)
    \item $\phi_{len}$: Perturb response length (short/normal/long)
    \item $\phi_{sent}$: Perturb response sentiment (negative/neutral/positive)
\end{itemize}
\end{definition}

\begin{definition}[Individual Bias Effect]
For a perturbation operator $\phi$, the individual bias effect is:
\begin{equation}
B_\phi = |S_\mathcal{M}(\phi(p_0)) - S_\mathcal{M}(p_0)|
\end{equation}
This measures how much a single perturbation changes the score.
\end{definition}

\begin{definition}[Multi-Perturbation Effect]
For a sequence of perturbation operators $\phi_1, \phi_2, ..., \phi_k$, the combined effect is:
\begin{equation}
C_{\phi_1,...,\phi_k} = |S_\mathcal{M}(\phi_1(\phi_2(...\phi_k(p_0)...))) - S_\mathcal{M}(p_0)|
\end{equation}
\end{definition}

\section{Interaction Ratio}

The central question of our work is: how does $C_{\phi_1,...,\phi_k}$ relate to the sum of individual $B_{\phi_i}$?

\begin{definition}[Interaction Ratio]
For a set of perturbation operators $\Phi = \{\phi_1, ..., \phi_k\}$, the Interaction Ratio (IR) is:
\begin{equation}
IR_\Phi = \frac{C_\Phi}{\sum_{i=1}^k B_{\phi_i}}
\end{equation}
where $C_\Phi$ is the combined effect and $B_{\phi_i}$ are the individual effects.
\end{definition}

\begin{definition}[Interaction Classification]
Given threshold $\epsilon > 0$ (typically $\epsilon = 0.05$):
\begin{itemize}
    \item If $IR_\Phi > 1 + \epsilon$: \textbf{Compounding} interaction
    \item If $1 - \epsilon \leq IR_\Phi \leq 1 + \epsilon$: \textbf{Additive} interaction
    \item If $IR_\Phi < 1 - \epsilon$: \textbf{Cancelling} interaction
\end{itemize}
\end{definition}

\section{Theorems}

\begin{theorem}[Additive Baseline]
If the score function $S_\mathcal{M}$ is linear in prompt perturbations, then $IR_\Phi = 1$ for all $\Phi$, i.e., all interactions are additive.
\end{theorem}

\begin{proof}
By linearity, $S_\mathcal{M}(\phi_1(\phi_2(p_0))) = S_\mathcal{M}(\phi_1(p_0)) + S_\mathcal{M}(\phi_2(p_0)) - S_\mathcal{M}(p_0)$. Therefore $C_{\phi_1,\phi_2} = |S_\mathcal{M}(\phi_1(p_0)) - S_\mathcal{M}(p_0) + S_\mathcal{M}(\phi_2(p_0)) - S_\mathcal{M}(p_0)| \leq B_{\phi_1} + B_{\phi_2}$. Under the assumption of aligned perturbation directions, $C_{\phi_1,\phi_2} = B_{\phi_1} + B_{\phi_2}$, giving $IR = 1$.
\end{proof}

\begin{theorem}[Model-Specific Interaction]
The Interaction Ratio $IR_\Phi$ is model-specific and depends on the internal representation of prompt features within model $\mathcal{M}$.
\end{theorem}

\begin{proof}[Sketch]
Consider two models $\mathcal{M}_1$ and $\mathcal{M}_2$ that differ in their training. Let $\mathcal{M}_1$ have independent feature representations for position and length, and $\mathcal{M}_2$ have shared representations. Under perturbation $\Phi = \{\phi_{pos}, \phi_{len}\}$:
- For $\mathcal{M}_1$: $IR \approx 1$ (independent features → additive)
- For $\mathcal{M}_2$: $IR \neq 1$ (shared features → interaction)

The existence of models with $IR \neq 1$ is empirically verified in our experiments.
\end{proof}

\begin{theorem}[Bias Amplification Bound]
For any model $\mathcal{M}$ and two perturbations $\phi_1, \phi_2$:
\begin{equation}
0 \leq C_{\phi_1,\phi_2} \leq B_{\phi_1} + B_{\phi_2} + 2\sqrt{B_{\phi_1}B_{\phi_2}}
\end{equation}
The upper bound corresponds to perfect compounding ($IR = 1 + \frac{2\sqrt{B_{\phi_1}B_{\phi_2}}}{B_{\phi_1}+B_{\phi_2}}$), and the lower bound to perfect cancelling ($IR = 0$).
\end{theorem}

\begin{proof}
The combined score change is $C = |f(p_0 + \delta_1 + \delta_2) - f(p_0)|$ where $\delta_i$ represent perturbations. By Taylor expansion:
\begin{equation}
C \approx |\nabla f^T(\delta_1+\delta_2) + \frac{1}{2}(\delta_1+\delta_2)^T H_f (\delta_1+\delta_2)|
\end{equation}
where $H_f$ is the Hessian. The cross term $\delta_1^T H_f \delta_2$ represents the interaction, bounded by $2\sqrt{B_{\phi_1}B_{\phi_2}}$ via Cauchy-Schwarz.
\end{proof}

\section{Empirical Framework}

\begin{definition}[Full-Factorial Design]
A $2 \times 3 \times 3$ full-factorial design with factors:
\begin{itemize}
    \item Position: $X_1 \in \{\text{first}, \text{second}\}$
    \item Length: $X_2 \in \{\text{short}, \text{normal}, \text{long}\}$
    \item Sentiment: $X_3 \in \{\text{negative}, \text{neutral}, \text{positive}\}$
\end{itemize}
The score model is:
\begin{equation}
S = \mu + \alpha_1 X_1 + \alpha_2 X_2 + \alpha_3 X_3 + \beta_{12} X_1 X_2 + \beta_{13} X_1 X_3 + \beta_{23} X_2 X_3 + \gamma_{123} X_1 X_2 X_3 + \epsilon
\end{equation}
where $\mu$ is the grand mean, $\alpha_i$ are main effects, $\beta_{ij}$ are two-way interactions, $\gamma_{123}$ is the three-way interaction, and $\epsilon$ is random error.
\end{definition}

\begin{definition}[Marginal Interaction Ratio]
For two perturbation types $i$ and $j$, the Marginal Interaction Ratio is:
\begin{equation}
MIR_{ij} = \frac{|C_{i,j}|}{|B_i| + |B_j|}
\end{equation}
where $B_i$ and $B_j$ are the marginal main effects and $C_{i,j}$ is the combined effect when both perturbations are applied simultaneously.
\end{definition}

\begin{definition}[Bayesian Interaction Posterior]
Given observed scores, the posterior distribution of the interaction ratio is:
\begin{equation}
P(IR_{ij} \mid \text{data}) \propto P(\text{data} \mid IR_{ij}) \cdot P(IR_{ij})
\end{equation}
where $P(IR_{ij})$ is a prior (typically log-normal centered at 1) and the likelihood is based on the observed score differences.
\end{definition}

\section{The Bias Composition Hypothesis}

\begin{definition}[Bias Composition Hypothesis]
The Bias Composition Hypothesis (BCH) states that the total bias in an LLM judge score can be decomposed into:
\begin{equation}
B_{total} = \sum_i B_i + \sum_{i<j} I_{ij} + \sum_{i<j<k} I_{ijk} + ...
\end{equation}
where $B_i$ are individual bias effects, $I_{ij}$ are pairwise interactions, and higher-order terms are interactions of three or more biases. Our work provides the first empirical test of this hypothesis.
\end{definition}

\begin{corollary}[Implication for Bias Mitigation]
If $IR_{ij} > 1$ for a pair of biases, then mitigating only one bias is insufficient. The remaining bias's effect is amplified by the presence of the unmitigated bias. Conversely, if $IR_{ij} < 1$, biases partially compensate for each other.
\end{corollary}

\begin{lemma}[Direction-Dependent Interaction]
The Interaction Ratio depends on the direction of each perturbation:
\begin{equation}
IR_{\phi_1,\phi_2}(d_1, d_2) \neq IR_{\phi_1,\phi_2}(-d_1, -d_2)
\end{equation}
in general, where $d_i \in \{-1, 0, 1\}$ indicates the perturbation direction.
\end{lemma}

\begin{proof}
This follows from the asymmetry of the Hessian $H_f$. For language models, perturbations in opposite directions may activate different attention patterns, leading to asymmetric interactions. Empirically, we observe this in Claude's unique verbosity preference (prefers concise vs. prefers verbose).
\end{proof}

\section{Generalized Bias Interaction Model}

We propose a generalized model for predicting bias interaction effects:

\begin{equation}
IR_{ij} = 1 + \theta_{ij} \cdot \rho_{ij} \cdot \sqrt{\frac{B_i \cdot B_j}{B_i + B_j}}
\end{equation}

where:
\begin{itemize}
    \item $\theta_{ij} \in [-1, 1]$: Model-specific interaction coefficient (learned from data)
    \item $\rho_{ij} \in [0, 1]$: Feature overlap between bias types $i$ and $j$ (from mechanistic analysis)
    \item $B_i, B_j$: Individual bias effect sizes (measured empirically)
\end{itemize}

This model predicts:
\begin{itemize}
    \item $\theta_{ij} > 0$: Compounding (shared mechanisms amplify effects)
    \item $\theta_{ij} = 0$: Additive (independent mechanisms)
    \item $\theta_{ij} < 0$: Cancelling (competing mechanisms)
\end{itemize}

\section{Conclusion}

We have presented a formal mathematical framework for understanding bias interaction effects in LLM-as-a-Judge. The framework provides:
\begin{enumerate}
    \item Precise definitions for bias effects and interactions
    \item The Interaction Ratio as a quantifiable metric
    \item Theorems bounding possible interaction patterns
    \item A Bayesian approach for quantifying uncertainty
    \item The Bias Composition Hypothesis for decomposing total bias
    \item A generalized model for predicting interactions from individual bias measurements
\end{enumerate}

This framework unifies the existing literature and provides a foundation for future research on bias interactions, mitigation strategies, and mechanistic understanding of bias in LLM judges.

\end{document}
